\newpage
\section*{Συμπεράσματα και μελλοντικές επεκτάσεις}

Η παρούσα εργασία παρουσίασε τη μετάβαση μιας ήδη \textit{containerized} εφαρμογής διαχείρισης αποθήκης από τοπική εγκατάσταση σε πλήρως αυτοματοποιημένη, \textit{cloud-native} αρχιτεκτονική στο \textit{Microsoft Azure}. Η υλοποίηση κάλυψε το σύνολο του κύκλου ζωής μιας σύγχρονης \textit{cloud} εφαρμογής: από τη δημιουργία της υποδομής μέσω \textit{Infrastructure as Code} (\textit{OpenTofu}) και την αυτοματοποιημένη ρύθμισή της (\textit{Ansible}), έως την ασφαλή διαχείριση μυστικών (\textit{HashiCorp Vault}), τις πρακτικές \textit{GitOps} (\textit{Jenkins}, \textit{ArgoCD}), την ενσωμάτωση σχεδιαστικών προτύπων ανθεκτικότητας στον κώδικα, τη \textit{serverless} εκτέλεση συγκεκριμένων λειτουργιών (\textit{Knative}), και την τεκμηρίωση επιπέδων ποιότητας υπηρεσίας μέσω πειραματικών δοκιμών φόρτου.

Η εμπειρία της υλοποίησης ανέδειξε ότι η πολυπλοκότητα μιας \textit{cloud-native} αρχιτεκτονικής δεν έγκειται μόνο στην αρχική εγκατάσταση κάθε επιμέρους τεχνολογίας, αλλά κυρίως στις αλληλεπιδράσεις μεταξύ τους --- όπως αναδείχθηκε χαρακτηριστικά από το περιστατικό \textit{infinite loop} στη \textit{serverless} συνάρτηση (Κεφάλαιο 6), όπου μια φαινομενικά απλή ενέργεια (επισήμανση αρχείου) είχε απρόβλεπτη επίδραση σε ολόκληρο τον μηχανισμό ενεργοποίησης γεγονότων. Παρά την ευρεία κάλυψη των τεχνολογιών και πρακτικών που ζητούνταν, η εργασία έχει συγκεκριμένους περιορισμούς, οι οποίοι αποτελούν φυσικές κατευθύνσεις μελλοντικής επέκτασης:

\begin{itemize}
  \item \textbf{Διαχείριση secrets του \textit{RabbitMQ}}: τα διαπιστευτήρια του \textit{RabbitMQ} καταγράφονται στο \textit{Vault} για λόγους συνέπειας, αλλά η πραγματική σύνδεση παραμένει βασισμένη σε μεταβλητές περιβάλλοντος, λόγω τεκμηριωμένου περιορισμού του επίσημου \textit{Docker image} στη διαχείριση διαπιστευτηρίων μέσω μεταβλητών κατά την εκκίνηση σε πρόσφατες εκδόσεις.
  \item \textbf{Αυτόματη ενημέρωση εικόνας (\textit{ArgoCD Image Updater})}: εξετάστηκε η χρήση αυτόματου μηχανισμού ενημέρωσης \textit{image tags} χωρίς ανθρώπινη παρέμβαση, αλλά η προσέγγιση αυτή απαιτεί πηγή τύπου \textit{Kustomize} ή \textit{Helm} και εισάγει επιπλέον αρχεία υπερκάλυψης (\textit{override files}) που περιπλέκουν την αναγνωσιμότητα του \textit{repository} χωρίς αντίστοιχο όφελος στην παρούσα κλίμακα· αντί αυτού, υλοποιήθηκε αυτόματη ενημέρωση του \textit{manifest} απευθείας από το \textit{Jenkins pipeline}.
  \item \textbf{Πρότυπο \textit{Circuit Breaker}}: δεν υλοποιήθηκε στην παρούσα φάση· θα αποτελούσε φυσική επέκταση του υπάρχοντος μηχανισμού \textit{Retry}, προστατεύοντας το σύστημα από αλυσιδωτές αστοχίες σε περίπτωση παρατεταμένης μη διαθεσιμότητας εξωτερικής υπηρεσίας.
  \item \textbf{Παρακολούθηση κόστους ανά πόρο (\textit{FinOps})}: η εκτίμηση κόστους μέσω \textit{Infracost} καλύπτει το επίπεδο της υποδομής (\textit{OpenTofu}), αλλά δεν επεκτείνεται σε ανάλυση κόστους ανά μεμονωμένο \textit{pod} σε πραγματικό χρόνο --- λειτουργία που θα απαιτούσε πρόσθετο εργαλείο (π.χ. \textit{OpenCost}), εκτός του πλαισίου της παρούσας εργασίας.
\end{itemize}

Συνολικά, η εργασία απέδειξε στην πράξη ότι η μετάβαση σε \textit{cloud-native} αρχιτεκτονική δεν είναι απλώς μεταφορά κώδικα σε διαφορετική υποδομή, αλλά υιοθέτηση ενός ολοκληρωμένου συνόλου πρακτικών --- αυτοματοποίησης, ασφάλειας, παρατηρητικότητας, και σχεδιαστικής ανθεκτικότητας --- που καθιστούν ένα σύστημα έτοιμο να λειτουργήσει αξιόπιστα σε πραγματικές συνθήκες παραγωγής.